% ═══════════════════════════════════════════════════════════════════
% ПРИЛОЖЕНИЯ E–G: ПОЛНЫЕ ТАБЛИЦЫ, ОБОЗНАЧЕНИЯ, ГЛОССАРИЙ
% ═══════════════════════════════════════════════════════════════════

\section{Полные таблицы данных}\label{app:fulldata}

\subsection{Семейство фаз $\delta=\pi/N$ и торможения для $N=2,\dots,14$}

\begin{center}
\small
\begin{tabular}{@{}c c c c c c@{}}
\toprule
$N$ & $\delta=\pi/N$ & $\delta^2/2$ & $\gamma=\delta^4/k{=}1$ &
$\delta_{\mathrm{eff}}$ & $\gamma$ ($k=22$) \\
\midrule
2  & 1.5707963 & 1.2337006 & 6.0880682 & 9.5631151 & 0.2767304 \\
3  & 1.0471976 & 0.5483115 & 1.2025814 & 1.2593403 & 0.0546628 \\
4  & 0.7853982 & 0.3084251 & 0.3805043 & 0.2988473 & 0.0172956 \\
5  & 0.6283185 & 0.1973921 & 0.1558545 & 0.0979263 & 0.0070843 \\
6  & 0.5235988 & 0.1370778 & 0.0751642 & 0.0393464 & 0.0034166 \\
7  & 0.4487990 & 0.1007102 & 0.0404989 & 0.0181767 & 0.0018409 \\
8  & 0.3926991 & 0.0771065 & 0.0237754 & 0.0093371 & 0.0010807 \\
9  & 0.3490659 & 0.0609266 & 0.0148499 & 0.0051835 & 0.0006750 \\
10 & 0.3141593 & 0.0493407 & 0.0097396 & 0.0030599 & 0.0004427 \\
11 & 0.2855993 & 0.0407771 & 0.0066572 & 0.0019009 & 0.0003026 \\
12 & 0.2617994 & 0.0342697 & 0.0047089 & 0.0012326 & 0.0002140 \\
13 & 0.2416609 & 0.0292000 & 0.0034119 & 0.0008245 & 0.0001551 \\
14 & 0.2243995 & 0.0251775 & 0.0025351 & 0.0005690 & 0.0001152 \\
\bottomrule
\end{tabular}
\end{center}

\subsection{Числа бинарного кода и сверка сертификата E}

\begin{center}
\small
\begin{tabular}{@{}l l l@{}}
\toprule
Величина & Значение & Сверка \\
\midrule
Стартовая точка $p_0$ & $(36,36)$ & конвейер тора \\
Фаза $\delta_T$ & $\pi/4$ & семейство $\pi/N$ \\
Посещено ячеек & $48$ & $t^*=48$ (теорема~\ref{th:E}) \\
События трения & $212$ & поток E: совпадение \\
Краевые ячейки & $432$ & поток E: совпадение \\
Длина цепного кода & $114$ & поток E: совпадение \\
Циклический сдвиг & инвариантен & код Фримена \\
Сумма кода mod $n$ & $88$ & лог протокола \\
Замыкание $\sum dx,\sum dy$ & $(0,0)$ & тор замкнут \\
\bottomrule
\end{tabular}
\end{center}

\subsection{Инварианты всех ступеней}

\begin{center}
\small
\begin{tabular}{@{}l l l l@{}}
\toprule
Ступень & $\rho$ / ранг & Сигнатура & Определитель \\
\midrule
Тор & $1$ ($H^1$) & $(0,0)$ & $1$ (унимодулярна) \\
K3 (полная) & $22$ ($b_2$) & $(3,19)$ & $1$ (унимодулярна) \\
K3 (прямые $+h$) & $20$ & $(1,19)$ & $64=8^2$ \\
Клейн $\Jac$ & $3$ & $(3,0)$ & $1$ (главная поляризация) \\
Ферма $N=15$ & $91$ & $(91,0)$ & $1$ (главная поляризация) \\
Ферма $N=30$ & $406$ & $(406,0)$ & $1$ (главная поляризация) \\
\bottomrule
\end{tabular}
\end{center}

\subsection{Спектральные дроби сертификата G}

\begin{center}
\small
\begin{tabular}{@{}l l l@{}}
\toprule
Решётка & $\lambda_1$ & Примечание \\
\midrule
$\ZZ[i]$ & $4\pi^2/1$ & тор, $\tau=i$ \\
$\ZZ[(1+\sqrt{-7})/2]$ & $4\pi^2/(\sqrt7/2)^2=16\pi^2/7$ & Клейн \\
$\ZZ[\sqrt{-3}]$ & $4\pi^2/3$ & $d=3$ \\
$\ZZ[\sqrt{-15}]$ & $4\pi^2/(\im\tau)^2$, $\pi/15$ единица &
  стенд \\
$\ZZ[\sqrt{-30}]$ & $4\pi^2/(\im\tau)^2$, $\pi/30$ единица &
  стенд \\
\bottomrule
\end{tabular}
\end{center}

\subsection{Слой Коула--Селберга (сертификат G4)}

\begin{center}
\small
\begin{tabular}{@{}l l@{}}
\toprule
Поле & Формула $\omega^2$ \\
\midrule
$d=4$ & $\Ga(1/4)^4/(16\pi)$ \\
$d=7$ & $\tfrac{21+\sqrt{-7}}{896}
  \bigl[\Ga(\tfrac17)\Ga(\tfrac27)\Ga(\tfrac47)\bigr]^2/\pi^2$ \\
$d=3$ & $\tfrac{3-\sqrt{-3}}{8}\Ga(\tfrac13)^6/(\pi^2 2^{8/3})$ \\
$d=15$ & $j$-пара в $\QQ(\sqrt5)$;
  $R=\tfrac{(2\sqrt5-3)+(\sqrt5-2)\sqrt{-3}}2$ \\
K3 & $\Ga(1/4)^4/(16\pi\sqrt2)$; индекс $\sqrt2/32$ \\
\bottomrule
\end{tabular}
\end{center}

\subsection{Показатели $\Ga$-произведения объёма (H2)}

Для $N=15$ показатели $c_k$ при $k=1,\dots,14$:
\[
  4,4,6,4,4,6,4,4,6,4,4,6,4,4;
\]
сумма $c_k=68$. Формула:
$\mathrm{vol}_h=(2\pi)^{32}\prod_k\Ga(k/15)^{-c_k}=1125$,
$\det V^2=1265625/256$, остаток тождества $1.01\cdot10^{-119}$.
Для $N=30$ показатели ненулевые при чётных
$k=2,4,\dots,28$ с тем же рисунком $(4,4,6)$ по циклам ---
согласовано с редукцией $\Ga(k/30)$ к уровням $15$ по чётным $k$;
нечётные $k$ образуют самостоятельную ступень (H3).

\subsection{Состав выборки V2 (сертификат B)}

Выборка характеров детерминирована: до трёх характеров на проводник
(первый, средний, последний в списке). Для $N=15$: $1+3+3$... по
проводникам $h_3=1\to1$, $h_5=6\to3$, $h_{15}=84\to3$; итого
$7$? Протокол фиксирует $21$ тест --- три данных обхода $(r,s)$
на каждый характер: $7\cdot3=21$. Для $N=30$: $1+3+3+3+3+3=16$
характеров $\times3$ обхода $=48$ тестов. Всего $69$ тестов;
максимум относительного отклонения $3.9\cdot10^{-36}$.

\section{Обозначения}\label{app:notation}

\begin{center}
\small
\begin{tabular}{@{}l l@{}}
\toprule
Обозначение & Смысл \\
\midrule
$(n,B,\lambda_0)$ & целочисленная тройка: порядок, индекс, порог \\
$\delta=\pi/N$ & спиновая фаза \\
$\gamma=\delta^4/k$ & торможение \\
$\delta_{\mathrm{eff}}=\delta^5/k$ & эффективная фаза \\
$\Delta_{Ch}$ & объединённая формула поправок \\
$b_{Ch}(n)$ & $1-\cos(2\pi/n)$ \\
$C_N$ & кривая Ферма $x^N+y^N=z^N$ \\
$g$ & род; $g=(N-1)(N-2)/2$ \\
$T_N$ & набор характеров $1\le a,b$, $a+b\le N-1$ \\
$d$ & проводник: $N/\gcd(N,a,b)$ \\
$h_d$ & число характеров проводника $d$ \\
$\eta_{i,j}$ & базисная голоморфная форма \\
$\Omega_{a,b}$ & $\Bfun(a/N,b/N)$ \\
$P_{a,b}(r,s)$ & период; $\tfrac1N\zeta^{ra+sb}\Omega_{a,b}$ \\
$J$ & форма пересечений (поляризация) \\
$\Lambda_d$ & решётка блока проводника $d$ \\
$\Ind(\Lambda_d)$ & индекс блока в насыщении \\
$Q$ & граница знаменателей (сертификат F) \\
$t^*$ & время завершения потока (сертификат E) \\
SNF & нормальная форма Смита \\
$dps$ & число точных цифр \texttt{mpmath} \\
\bottomrule
\end{tabular}
\end{center}

\section{Глоссарий}\label{app:glossary}

\begin{description}[leftmargin=1.6em, itemsep=0.2em]
\item[Стенд] геометрический объект, на котором конвейер программы
  запускается целиком; ступень лестницы.
\item[Калибровочный стенд] нижняя ступень с полностью явным выводом
  (тор); служит эталоном для верхних ступеней.
\item[Конвейер] последовательность правил
  \eqref{eq:delta}--\eqref{eq:deltaCh}, переводящая тройку в полную
  структуру поправок.
\item[Сертификат] независимая проверка конкретного вычисления
  (точная арифметика либо прямое интегрирование).
\item[Замкнутая форма] выражение величины через $\Ga$-значения и
  корни единицы без численных подгонок.
\item[Проводник] наименьший уровень $d$, на котором характер
  определён; $d=N/\gcd(N,a,b)$.
\item[Изотипная компонента] собственное подпространство действия
  группы, отвечающее характеру.
\item[Насыщение] пересечение рационального расширения подрешётки с
  объемлющей решёткой.
\item[Errata] зафиксированное исправление с автоматической
  проверкой (E8 $28/36$, кейс E6).
\item[Плиточная система] модульная сетка диаграмм высокого
  разрешения в отчётах лаборатории.
\item[Лестница] последовательность стендов возрастающего масштаба с
  общими правилами конвейера.
\item[Тройка] $(n,B,\lambda_0)$ --- всё, что геометрия поставляет
  конвейеру.
\end{description}
